\section{The main theorem}
\label{sec:mai}

This section assembles Theorem~A and derives Corollary~B; the consequences of $\Dn=1$ for the natural form and universality are in \cref{sec:cmp}. A conifold frame of \cref{def:set-frames}, given by $(\alpha,\beta)\in\bbQ^2$, is fixed, and $M=\Mo(\tau(u)\zh)$ is taken exactly in $u$. As in \cref{sec:flh}, $\Lambda_s\coloneqq\diag(s^{1-\nu_a})$ and $\Lambda_\varsigma\coloneqq\diag(\varsigma^{1-\nu_a})$ act on lowered components (local letters, not the leg map of \cref{def:mod-weights}), and $1_{\TFT}=\sum_i\Delta_i^{-1/2}\tilde e_i$.

\subsection{Theorem A and Corollary B}
\label{sec:mai-thm}

\begin{thm}[Theorem A]\genusstatus{AI-proposed, unchecked}
\label{thm:mai-main}
Let $g\geq2$, let $(\alpha,\beta)\in\bbQ^2$ give a conifold frame (\cref{def:set-frames}), let $\varpi_1=u+O(u^2)$ be the vanishing period and $\tc=\varpi_1/\varpi_0$. Let $M\coloneqq\Mo(\tau(u)\zh)$, exactly in $u$, and $R'\coloneqq\Qr M^{-1}$, where $\Qr$ is the matrix of the lowered rows of \cref{def:set-rows} expanded in the frame (\cref{def:set-confexp}). For a leg-free stable graph $\Gamma$ of genus $g$ let $T_\Gamma\coloneqq\mathrm{Cont}_\Gamma[E_{R'},q_{R'};M.\TFT]$ be its outer term, $\TFT$ the TFT of \cref{def:set-frob} expanded in the frame, $\TFT_{g,n}(\tilde e_i,\dots,\tilde e_i)=(\zeta^{4i}s\varpi_0)^{2g-2+n}$ (\cref{def:set-outer}). Then:
\begin{enumerate}
\item $T_\Gamma$ lies in $\bbQ\fps{u}$ for every $\Gamma$ with at least one edge;
\item for the vertex graph, $T_\Gamma\in c_g\,\tc^{2-2g}+\bbQ\fps{u}$, with
\[
c_g=\frac{(-1)^{g-1}B_{2g}\,5^{1-g}}{2g(2g-2)}
\]
as in \eqref{eq:set-cg}.
\end{enumerate}
The hatted terms are $G_\Gamma=u^{2g-2}T_\Gamma$ (\cref{lem:flh-hat}).
\end{thm}

A priori $T_\Gamma\in\bbQ(\zeta)\lser{s}$; the vertex graph is that of \cref{def:set-outer}.

\begin{proof}
Fix $g\geq2$ and the frame. By \cref{thm:gam-D}, $\Dn_0=1$ in $\bbQ\fps{z}$; in particular $\Fg{g}$ holds (\cref{def:gau-Fg}). By the comparison theorem, \cref{thm:flh-main}, $T_\Gamma=F_\Gamma$ for every leg-free stable graph $\Gamma$ of genus $g$, where $F_\Gamma$ is the term of $\Gamma$ in the flat sum of $\Pc.\Omega^{\mathrm{mod}}_{v(u)}$ (\cref{def:red-flatsum}). By \cref{thm:mod-MLK} the Model Lemma $\ML{h}$ holds for every $h\geq0$. So \cref{cor:red-can} applies: since $\ML{h}$ holds for $h\leq g-1$, $F_\Gamma$ lies in $\bbQ\fps{u}$ for every $\Gamma$ with an edge; since $\ML{g}$ holds as well, the vertex graph gives $F_\Gamma\in c_g\tc^{2-2g}+\bbQ\fps{u}$, the leaf point of $\Pc$ having the vanishing coordinate $T_{\Pc}=\tc$ (\cref{prop:tra-can}). This proves (i) and (ii). The last sentence is \cref{lem:flh-hat}.

The statements cited rest on these: \cref{thm:gam-D} on \cref{thm:con-glue}, Theorem~P (\cref{thm:con-Pa,thm:con-Pb}), \cref{thm:gam-main}~(ii), \cref{prop:gam-hatflat}, \cref{cor:gau-colours}; \cref{thm:flh-main} on \cref{thm:gau-R0}~(iv), \cref{cor:gau-colours}, \cref{prop:gau-symp}, \cref{lem:set-M0sympl}; \cref{prop:tra-can} on \cref{thm:gau-U,thm:gau-G,prop:gau-symp} and \cref{sec:tra}; all on \cref{prop:gau-quintic}, from \cref{prop:set-cyclicLR,thm:gau-cyc}.
\end{proof}

\emph{Where the inputs enter.} (QRR) enters through the rows: its flatness part (\cref{def:set-rows}~(b)) in \cref{prop:set-cyclicLR}~(iii) and \cref{cor:set-cyclicPF5}~(b), its value (\cref{thm:set-qrr}) in \cref{lem:gau-colrows} and the glue \cref{thm:con-glue}; (M) and (G0) through \cref{thm:set-mirror,prop:set-cyclicLR,cor:set-cyclicPF5}; (T) through \cref{thm:set-GT}, which identifies the $R$-matrix whose rows define $T_\Gamma$; \cite{PPZ15} through \cref{thm:set-G0} (as part of (G0)), \cref{def:set-graphsum}, \cref{lem:set-pushins} (in the reading of the topological-recursion input, \cref{prop:mod-Rtr}~(ii)), \cref{def:set-modcohft}, \cref{lem:set-modhom}, \cref{lem:set-dilaton}, \cref{def:set-outer}, \cref{sec:tra-dress}, \cref{thm:mod-tr} and \cref{sec:flh}; \cite{Eynard:2007kz}, \cite{DuninBarkowski:2012bw} (as recalled in \cite{Dunin-Barkowski:2015caa}) and \cite{ABDKS25} through \cref{thm:mod-MLK}, on the model curve $x=\Wv(p)$, $y=p$ only; and \cite[(1.7)]{AC96} through \cref{lem:set-dilaton}.

\emph{Conditionality.} Theorem~A holds modulo the quoted theorems (Q) of \cite{Eynard:2007kz}, \cite{DuninBarkowski:2012bw} (as recalled in \cite{Dunin-Barkowski:2015caa}), \cite{ABDKS25} and \cite{PPZ15} with the statement numbers listed in \cref{sec:int-inputs}, the flatness of the $R$-matrix and its value at the large-radius point (QRR), the genus-$0$ mirror theorem (M) and the genus-$0$ axioms (G0) of the twisted theory, Teleman's reconstruction with the axioms of the twisted theory in all genera (T), and the formula \cite[(1.7)]{AC96}; every other step is proved in the note. Topological recursion is used only on the model curve, never on a mirror curve of the quintic or of the formal quintic, and no finite computation enters the proof.

To expand $F_g$ (\cref{def:set-fq}) at the conifold point we use its graph-sum representative in $\Nring$ (\cref{def:set-confexp}). Write $\Psi_{\mathrm Q}\coloneqq\Qr|_{z=0}$, now for the rows as elements of $\Nring\otimes\bbQ(\zeta)$ (\cref{cor:gau-rowsNN}); its conifold expansion $\Psi_{\mathrm{cf}}\coloneqq\Cexp_{\alpha,\beta}(\Psi_{\mathrm Q})$ is the matrix $\Psi_{\mathrm Q}$ of \cref{lem:flh-inputs}, and $\Psi(u)\coloneqq\Lambda_s\Psi_{\mathrm{cf}}$ is its hatted form; $\mathfrak m_a$, $\mathfrak w_a$ are local letters.

\begin{lem}[the matrices $\Psi_{\mathrm Q}$ and $\Psi(u)$]\genusstatus{AI-proposed, unchecked}
\label{lem:mai-Psi}
\begin{enumerate}
\item In $\Nring\otimes\bbQ(\zeta)$, $(\Psi_{\mathrm Q})_{aj}=\zeta^{j(a+1)}\mathfrak m_a$, where
\[
(\mathfrak m_0,\dots,\mathfrak m_4)=\bigl(L/C_0,\;L^2/(C_0C_1),\;C_0C_1/L^2,\;C_0/L,\;1\bigr),
\]
and $\Psi_{\mathrm Q}^T\eta^{-1}\Psi_{\mathrm Q}=1$.
\item $\Psi(u)_{aj}=\zeta^{j(a+1)}\mathfrak w_a(u)$, where
\[
(\mathfrak w_0,\dots,\mathfrak w_4)=\bigl(1/\varpi_0,\;1/(\varpi_0\Dq\tc),\;\varpi_0\Dq\tc,\;\varpi_0,\;1\bigr).
\]
Hence $\Psi(u)\in GL_5(\bbQ(\zeta)\fps{u})$, $\Psi(0)=\Psi_0$, and $\Psi_{\mathrm{cf}}^T\eta^{-1}\Psi_{\mathrm{cf}}=\Psi(u)^T\eta^{-1}\Psi(u)=1$, equivalently $\Psi_{\mathrm{cf}}\Psi_{\mathrm{cf}}^T=\Psi(u)\Psi(u)^T=\eta$.
\end{enumerate}
\end{lem}

\begin{proof}
(i) At $z=0$, $\Qr^{(j)}_0(0)=\eta(1,\tilde e_j)=\Delta_j^{-1/2}=\zeta^jL/C_0$ (\cref{def:set-frob}); equivalently $\mathrm{ph}_0=1$ (\cref{def:set-rows}). By \cref{prop:set-cyclicLR}~(iii), $X_a=\zeta^{-j}\Qr^{(j)}_a$ satisfies $X_{a+1}=(C_1\kappa^{\mathrm{LR}}_a)^{-1}(z\Dq+\zeta^jL)X_a$ for $a\leq3$; at $z^0$ the term $z\Dq$ drops out, so $X_{a+1}(0)=\zeta^jL\,X_a(0)/(C_1\kappa^{\mathrm{LR}}_a)$, with $C_1\kappa^{\mathrm{LR}}=(C_1,L^5/(C_0^2C_1^2),C_1,C_0,C_0)$ and $X_0(0)=L/C_0$. This gives $\Qr^{(j)}_a(0)=\zeta^{j(a+1)}\mathfrak m_a$; only the relations of \cref{prop:set-cyclicLR} are used, so the identity holds in $\Nring\otimes\bbQ(\zeta)$ (\cref{cor:gau-rowsNN}). Since $(\eta^{-1})_{ab}=\frac15\delta_{a+b\equiv3}$ and $\mathfrak m_a\mathfrak m_{3-a}=1$ for every $a$,
\[
(\Psi_{\mathrm Q}^T\eta^{-1}\Psi_{\mathrm Q})_{ij}=\frac15\sum_a\zeta^{i(a+1)+j(4-a)}=\frac15\,\zeta^{i-j}\sum_a\zeta^{(i-j)a}=\delta_{ij}.
\]
(ii) The conifold expansion ($C_0\mapsto\varpi_0$, $C_1\mapsto\Dq\tc$, $L\mapsto s^{-1}$, \cref{lem:gau-confexp}) followed by the hatting factor $s^{1-\nu_a}$, $\nu=(0,-1,3,2,1)$, turns $\mathfrak m_a$ into $\mathfrak w_a$. At $u=0$, $\varpi_0(0)=1$ and $\Dq\tc(0)=-1$ (\cref{def:set-frames}) give $(\mathfrak w_a(0))_a=(1,-1,-1,1,1)=(\varepsilon_a)_a$, so $\Psi(0)=\Psi_0$ (\cref{def:set-M0}). The $\mathfrak w_a$ are units of $\bbQ\fps{u}$ with $\mathfrak w_a\mathfrak w_{3-a}=1$, which gives $\Psi(u)^T\eta^{-1}\Psi(u)=1$ as in (i); and $\Psi_{\mathrm{cf}}^T\eta^{-1}\Psi_{\mathrm{cf}}=\Psi(u)^T\Lambda_s^{-1}\eta^{-1}\Lambda_s^{-1}\Psi(u)=\Psi(u)^T\eta^{-1}\Psi(u)$, as $\Lambda_s\eta^{-1}\Lambda_s=\eta^{-1}$ ($\nu_a+\nu_{3-a}=2$). The equivalent forms follow from $A^{-1}=A^T\eta^{-1}$ for $A=\Psi_{\mathrm{cf}}$ and $A=\Psi(u)$.
\end{proof}

\begin{defn}[the genus-$g$ sum in the ring and its conifold expansion]
\label{def:mai-Fcf}
By \cref{cor:gau-rowsNN} the row coefficients $\Qr^{(i)}_{a,n}$ lie in $\Nring\otimes\bbQ(\zeta)$ and are free of $\logq$. By \cref{lem:mai-Psi}~(i), $\Psi_{\mathrm Q}^T\eta^{-1}\Psi_{\mathrm Q}=1$ in $\Nring\otimes\bbQ(\zeta)$. Put
\[
\tilde R\coloneqq\Psi_{\mathrm Q}^{-1}\Qr=\Psi_{\mathrm Q}^T\eta^{-1}\Qr ,
\]
the $R$-matrix in Teleman's colour-colour form; $\tilde R(0)=1$. For a matrix series $X(z)$ with $X(0)=1$ let $N_X(\psi,\psi')\coloneqq1-X^{-1}(\psi)X^{-1}(\psi')^T$ and
\[
E_X(\psi,\psi')\coloneqq\frac{N_X(\psi,\psi')-\frac12N_X(-\psi',\psi')-\frac12N_X(\psi,-\psi)}{\psi+\psi'} ,
\]
a power series for every $X$, since its numerator vanishes at $\psi=-\psi'$; it is symmetric, $E_X(\psi,\psi')^T=E_X(\psi',\psi)$, as $N_X(\psi,\psi')^T=N_X(\psi',\psi)$, and it is the edge of \cite[Def.~2.2]{PPZ15} when $X(-z)^TX(z)=1$, because then $N_X(-\psi',\psi')=N_X(\psi,-\psi)=0$ \cite[Rem.~2.1, (6)]{PPZ15}. (The symplecticity of $\tilde R$ is not known in $\Nring\otimes\bbQ(\zeta)$.) With the leaf point $q_{\tilde R}=\tilde R^{-1}(z)(-z1_{\TFT})$, $1_{\TFT}=\sum_i\Delta_i^{-1/2}\tilde e_i$, define
\[
F_g^{\Nring}\coloneqq\sum_\Gamma\mathrm{Cont}_\Gamma\bigl[E_{\tilde R},q_{\tilde R};\TFT\bigr]\in\Nring\otimes\bbQ(\zeta),
\]
the sum over the leg-free stable graphs of genus $g$ ($|\Aut\Gamma|^{-1}$ is part of $\mathrm{Cont}_\Gamma$). The \emph{conifold expansion} of $F_g$ in the frame $(\alpha,\beta)$ is
\[
F_g^{\mathrm{cf}}\coloneqq\Cexp_{\alpha,\beta}\bigl(F_g^{\Nring}\bigr)\in\bbQ\lser{s}.
\]
\end{defn}

The definition is justified as follows. First, $F_g^{\Nring}$ is a polynomial over $\bbQ$ in finitely many row coefficients, in $\Delta_i^{\pm1/2}=(\zeta^{4i}C_0/L)^{\pm1}$ and in $\zeta$: the contraction is finite because $q_{\tilde R}+z1_{\TFT}=z(1_{\TFT}-\tilde R^{-1}(z)1_{\TFT})$ is $O(z^2)$ (\cref{def:flh-outer}~(ii)), the coefficients of $\tilde R^{-1}$ are polynomials in those of $\tilde R$, and the correlators of $\TFT$ are rational multiples of $\Delta_i^{(2h-2+n)/2}$, $2h-2+n>0$. Second, $F_g^{\Nring}$ lies in $\Nring$. For $k\in(\bbZ/5)^\times$ let $\sigma_k$ be the automorphism $\zeta\mapsto\zeta^k$ of $\Nring\otimes\bbQ(\zeta)$. By \cref{cor:gau-rowsNN}~(iii), which rests on \cref{lem:gau-colrows}, $\Qr^{(i)}_{a,n}=\zeta^{i(a+1-n)}\Qr^{(0)}_{a,n}$; and the colour-$0$ representatives lie in $\Nring$: by \cref{thm:con-glue}, $\Dn_0\in1+z\bbQ\fps z$, so $A^{\mathrm{act}}_n=\sum_{m\leq n}A^{\mathrm{can}}_m\coef{z^{n-m}}\Dn_0^{-1}$ lies in $\bbQ[L,L^{-1}]$ (\cref{thm:gau-R0}~(ii), (iii)), and the representatives of \cref{cor:gau-rowsNN}~(i) are obtained from $A^{\mathrm{act}}_n(L)\,c_0^{-1}$ by a recursion with coefficients in $\Nring$. Moreover $\Delta_i^{1/2}=\zeta^{4i}C_0/L$; hence $\sigma_k(\Qr^{(i)}_{a,n})=\Qr^{(ki)}_{a,n}$, $\sigma_k(\Delta_i^{1/2})=\Delta_{ki}^{1/2}$, and, by the explicit form of $\Psi_{\mathrm Q}$ in \cref{lem:mai-Psi}, $\sigma_k(\tilde R_{ij})=\tilde R_{ki,kj}$. So $\sigma_k$ permutes the colours; as $\TFT$ is diagonal in the colour basis, $F_g^{\Nring}$ is a sum over all colourings of the vertices, which $\sigma_k$ permutes, so $\sigma_k(F_g^{\Nring})=F_g^{\Nring}$ for every $k$ and $F_g^{\Nring}\in(\Nring\otimes\bbQ(\zeta))^{\mathrm{Gal}}=\Nring$. Third, $\ev(F_g^{\Nring})=F_g$: under $\ev$, $\tilde R$ becomes the $R$-matrix of \cref{def:set-rows} in the normalised canonical basis, which is symplectic (\cref{def:set-rows}~(a)), so the image of $E_{\tilde R}$ is the edge of \cite[Def.~2.2]{PPZ15}; the graph sum is then the integral over $\overline{\mathcal M}_g$ of the genus-$g$ class of the unit-preserving action on $\TFT$ (\cref{lem:set-pushins}~(iii), unit $1_{\TFT}$), which is the potential by (T), with the normalisation (QRR) (\cref{thm:set-GT,thm:set-qrr}). A change of the primitive of $1$ in \cref{def:set-confexp} fixes the rows and $\Delta_i^{1/2}$, hence $F_g^{\Nring}$. Whether $\ev$ is injective on $\Nring$ is neither needed nor proved.

\begin{cor}[Corollary B: the conifold gap of the formal quintic]\genusstatus{AI-proposed, unchecked}
\label{cor:mai-gap}
Let $g\geq2$. In every conifold frame of \cref{def:set-frames},
\[
F_g^{\mathrm{cf}}\in c_g\,\tc^{2-2g}+\bbQ\fps{u}.
\]
Since $\tc=u+O(u^2)$ lies in $\bbQ\fps{u}$, $\bbQ\fps{u}=\bbQ\fps{\tc}$: the conifold expansion of $F_g$ contains no fractional power of $u$ and no power $\tc^k$ with $3-2g\leq k\leq-1$, and its leading term is $c_g\tc^{2-2g}$.
\end{cor}

\begin{proof}
Write $\Cexp$ for $\Cexp_{\alpha,\beta}$, extended $\bbQ(\zeta)$-linearly.

\emph{Step 1: expanding the graph sum.} By \cref{lem:gau-confexp}, $\Cexp$ is a ring homomorphism; $F_g^{\Nring}$ is a polynomial expression in finitely many coefficients of the edges of \cref{def:mai-Fcf}, of the leaf points and of $\TFT$, and the edge $E_X$, polynomial in the coefficients of $X^{-1}$ followed by an exact division by $\psi+\psi'$, commutes with ring homomorphisms. Hence
\[
F_g^{\mathrm{cf}}=\sum_\Gamma\mathrm{Cont}_\Gamma\bigl[E_{\Cexp(\tilde R)},q_{\Cexp(\tilde R)};\Cexp(\TFT)\bigr].
\]
Since $\Dn=1$ (\cref{thm:gam-D}), \cref{lem:flh-inputs}~(iii) gives $R'(-z)^T\eta^{-1}R'(z)=1$ for $R'=\Cexp(\Qr)M^{-1}$, and $M(-z)^TM(z)=1$ (\cref{lem:set-M0sympl}); so $\Cexp(\Qr)(-z)^T\eta^{-1}\Cexp(\Qr)(z)=1$. With $\Psi_{\mathrm{cf}}=\Cexp(\Psi_{\mathrm Q})$, so that $\Psi_{\mathrm{cf}}^T\eta^{-1}\Psi_{\mathrm{cf}}=1$ and $\Psi_{\mathrm{cf}}\Psi_{\mathrm{cf}}^T=\eta$ (\cref{lem:mai-Psi}), and $\Cexp(\tilde R)=\Psi_{\mathrm{cf}}^{-1}\Cexp(\Qr)$, this gives
\[
\Cexp(\tilde R)(-z)^T\,\Cexp(\tilde R)(z)=\Cexp(\Qr)(-z)^T(\Psi_{\mathrm{cf}}\Psi_{\mathrm{cf}}^T)^{-1}\Cexp(\Qr)(z)=1 :
\]
$\Cexp(\tilde R)$ is symplectic, and $E_{\Cexp(\tilde R)}$ is the edge of \cite[Def.~2.2]{PPZ15}. The classes of $\Cexp(\TFT)$ are those of $\TFT$ with $\Delta_i^{\pm1/2}$ replaced by $\Cexp(\Delta_i^{\pm1/2})$, and the axioms of a topological field theory with unit are polynomial identities in the $\Delta_i^{\pm1/2}$, so $\Cexp(\TFT)$ is one over $\bbQ(\zeta)\lser{s}$, with the unit $\Cexp(1_{\TFT})$, and $q_{\Cexp(\tilde R)}=\Cexp(\tilde R)^{-1}(z)\bigl(-z\,\Cexp(1_{\TFT})\bigr)$. By \cref{lem:set-pushins}~(iii), the right-hand side above is $\int_{\overline{\mathcal M}_g}\bigl(\Cexp(\tilde R).\Cexp(\TFT)\bigr)_{g,0}$, for the unit-preserving action of \cite[Def.~2.13]{PPZ15}.

\emph{Step 2: regrouping.} $M$ is symplectic with $M(0)=1$ (\cref{lem:set-M0sympl}), so $X\coloneqq\Cexp(\tilde R)M^{-1}$ is symplectic with $X(0)=1$, and $\Cexp(\Qr)M^{-1}=\Psi_{\mathrm{cf}}X$ is the outer matrix $R'$ of \cref{def:set-outer}. By \cite[Prop.~2.14]{PPZ15}, the unit-preserving action of \cite[Def.~2.13]{PPZ15} is a left group action, so
\[
\Cexp(\tilde R).\Cexp(\TFT)=X.\bigl(M.\Cexp(\TFT)\bigr).
\]
(The results of \cite{PPZ15} are stated over $\bbQ$; \cite[Props.~2.12 and~2.14]{PPZ15} hold over $\bbQ(\zeta)\lser s$ for the topological field theory $\Cexp(\TFT)$, whose unit has the components $y_i=\Cexp(\Delta_i^{-1/2})$, units of $\bbQ(\zeta)\lser s$, in the orthonormal basis $\tilde e_i$, and for the symplectic series $M$ and $X$, by the last paragraph of \cref{def:set-graphsum}.) Take the genus-$g$ parts without legs and integrate over $\overline{\mathcal M}_g$. The left-hand side gives $F_g^{\mathrm{cf}}$ by Step 1; the right-hand side gives $\sum_\Gamma\mathrm{Cont}_\Gamma[E_X,q_X;M.\Cexp(\TFT)]$ by \cref{lem:set-pushins}~(iii), applied to $M.\Cexp(\TFT)$, whose unit is that of $\Cexp(\TFT)$ by \cite[Prop.~2.12]{PPZ15}, and to the symplectic series $X$, $X(0)=1$; this is the contraction of \cref{def:flh-outer} with $\varkappa=1$. Since $\Psi_{\mathrm{cf}}$ is an isometry from the colour space onto $(\bbH,\eta)$ independent of $z$ which maps $1_{\TFT}$ to $1$ (\cref{lem:mai-Psi}~(ii), \cref{lem:flh-inputs}~(iv)), $R'(-\psi)^T\eta^{-1}R'(-\psi')=X(-\psi)^TX(-\psi')$ and $R'^{-1}(z)(-z1)=X^{-1}(z)(-z1_{\TFT})$; so $E_{R'}=E_X$ and $q_{R'}=q_X$: $R'^{-1}=X^{-1}\Psi_{\mathrm{cf}}^{-1}$, and $\Psi_{\mathrm{cf}}^{-1}\otimes\Psi_{\mathrm{cf}}^{-1}$ maps the inverse bivector of $(\bbH,\eta)$ to that of the colour space, $\Psi_{\mathrm{cf}}$ being an isometry. Each term is therefore $\mathrm{Cont}_\Gamma[E_{R'},q_{R'};M.\TFT]=T_\Gamma$, with $\TFT$ expanded in the frame as in \cref{def:set-outer}, and
\[
F_g^{\mathrm{cf}}=\sum_\Gamma T_\Gamma ,
\]
with no further factor $|\Aut\Gamma|^{-1}$.

\emph{Step 3.} By Theorem~A (\cref{thm:mai-main}), $T_\Gamma\in\bbQ\fps{u}$ for every $\Gamma$ with an edge, and the vertex graph gives $c_g\tc^{2-2g}+\bbQ\fps{u}$; the sum is finite, so $F_g^{\mathrm{cf}}\in c_g\tc^{2-2g}+\bbQ\fps{u}$.
\end{proof}

Corollary~B uses Theorem~A, the reconstruction (T) through the identification $\ev(F_g^{\Nring})=F_g$, \cite[\S0.2, (1); Def.~0.3, in particular (iii); (5); Rem.~2.1 and (6); Defs.~2.2 and~2.5, with the conventions stated after them; (8); Prop.~2.12; Def.~2.13; Prop.~2.14]{PPZ15}, and proved statements only; its conditionality is that of Theorem~A.

\begin{rmk}[the constant]
\label{rmk:mai-constant}
The constant $c_g$ is that of \cite[Conj.~1.5]{CGYZ26}, proved there for the quintic threefold \cite[Thm.~2]{CGYZ26}, in the same normalisation of the vanishing period; the two theories have different potentials, so nothing about the quintic follows. For the values see \cref{rmk:set-cg}; for the dictionary, the rescaling of the two periods of a frame, and the statement that $F_g$ lies in the algebra of \cite[Def.~1.1]{CGYZ26}, so that for an admissible pair $F_g^{\mathrm{cf}}$ is its conifold expansion in the sense of \cite[Def.~1.2]{CGYZ26}, see \cref{sec:sgn-lit}.
\draftnote{Numbers from arXiv:2609.13921v2, read in the HTML rendering; to be re-read in the PDF.}
\end{rmk}

\begin{rmk}[graph-locality]
\label{rmk:mai-graphlocal}
In the normalisation of \cref{rmk:mai-stored}, a computation that accompanied the work, not recorded in \cref{sec:evi}, found six edge graphs of genus $4$ singular while the singular part of the total was unchanged. If so, the singular part of the total does not determine those of the graph terms, and no argument that uses the total only through its singular part can prove Theorem~A; this is not used. The proof is graph-local, and it uses the odd part of $\log\cn$, which is fixed by (G) (\cref{thm:gau-G}) and not by symplecticity.
\end{rmk}

\begin{rmk}[the hypotheses cannot be dropped]
\label{rmk:mai-necessity}
Computations that accompanied the work, not recorded in \cref{sec:evi}, found singular edge graphs in the flat sums of gauges that violate (A), among them $P^\sharp$ of \cref{rmk:mai-stored}, and of gauges that are not regular in $u$ in the flat frame (and so are not dressings in the sense of \cref{def:tra-dressing}). This indicates, but does not prove here, that neither hypothesis of \cref{thm:red-main} can be dropped in general. For (Sp) no such example exists, because without it the edge of \cref{def:tra-edge} is not a power series. The gauge equation is not a hypothesis (\cref{rmk:mai-annih}).
\end{rmk}

